\documentclass{article}
\usepackage[utf8]{inputenc}
\usepackage[margin=2cm]{geometry}
\usepackage{fullpage}
\usepackage{enumitem}
\usepackage{amssymb}
\usepackage{amsmath}
\usepackage{cancel}
\usepackage{gensymb}
\usepackage{graphicx}
\usepackage{indentfirst}
\usepackage{listings}
\usepackage{algpseudocode}
\usepackage{amsmath}
\usepackage{float}
\usepackage{url}
\usepackage{tikz}
\usepackage{hyperref}
\usepackage{color}
\usepackage{upgreek}
\usepackage[russian]{babel}
\usepackage{amsthm}
\usepackage{ wasysym }
\usepackage{tipa}
\usepackage{todonotes}
\usepackage{mathrsfs}

\theoremstyle{definition}

\newtheorem{theorem}{Теорема}[section]
\newtheorem{conjecture}[theorem]{Утверждение}
\newtheorem{definition}[theorem]{Определение}
\newtheorem{lemma}[theorem]{Лемма}
\newtheorem{corollary}[theorem]{Следствие}
\newtheorem{example}[theorem]{Пример}

\newcommand{\An}{\operatorname{An}}
\newcommand{\G}{\mathcal{G}}
\newcommand{\pattern}{\mathcal{P}}
\newcommand{\Gsym}{\mathcal{G}_s}
\newcommand{\OPT}{\operatorname{OPT}}
\newcommand{\D}{\Delta}
\newcommand{\hGamma}{\hat{\Gamma}}


\author{Калинку Максим}
\title{Наибольшее сборное число для сборного графа \\ с заданным количеством вершин\\ Черновик}

\begin{document}
\maketitle

\section*{Аннотация}

Сборным числом простого сборного графа называется минимальный размер гамильтонова множества полигональных путей, покрывающего весь граф. Реализующим числом целого числа $n$ называется наименьший возможный размер графа, сборное число которого равно $n$. Недавние исследования показали, что известная верхняя оценка для реализующего числа является точной: в общем классе сборных графов $R_{\min}(n)=3n-2$, а множество графов, максимизирующих сборное число при заданном числе вершин, состоит в точности из петельных подстановок. Поскольку это семейство содержит много петель, в настоящей работе мы переносим внимание на класс сборных графов без петель. Получены новые верхние и нижние оценки для $R_{\min}^\circ$ в этом ограниченном семействе; доказано, что $R_{\min}^\circ(n)=5n-3$, и дана полная характеризация графов, на которых достигается максимум сборного числа при заданном размере в классе графов без петель.

\medskip
\noindent\textbf{Ключевые слова:} сборное число, два-слово, простой сборный граф, реализующее число, полигональный путь, гамильтоново множество.

\section{Введение}

Понятие простого сборного графа было введено в работе~\cite{Angeleska2009} как математическая модель процесса ДНК-рекомбинации у некоторых видов инфузорий. С точки зрения комбинаторики, эта модель представляет собой конечный связный граф с вершинами степени $1$ и $4$. В каждой вершине степени $4$ задан циклический порядок инцидентных полурёбер. Кроме того, в графе существует трансверсальный эйлеров путь. Под мощностью простого сборного графа $|\Gamma|$ понимается число четырёхвалентных вершин в нём.

Биологический смысл этой конструкции заключается в следующем: жёсткая вершина отвечает сайт-специфической рекомбинации, а отдельный ген, образующийся в результате рекомбинации, моделируется \emph{полигональным путём} — путём, который в каждой вершине степени 4 «поворачивает», то есть проходит через неё по соседним в циклическом порядке полурёбрам. Множество генов, полученных из одной макромолекулы, даёт \emph{гамильтоново множество полигональных путей} — набор попарно не пересекающихся по вершинам полигональных путей, покрывающий все $4$-валентные вершины графа~\cite{Muche2012}. Минимальное количество путей среди всех таких наборов называется \emph{сборным числом} графа $\Gamma$ и обозначается $\An(\Gamma)$.

С каждым простым сборным графом естественно связано \emph{два-слово} — слово, в котором каждый символ некоторого алфавита появляется ровно дважды. Вершины графа отождествляются с различными буквами, а последовательные пары букв два-слова задают рёбра эйлерова трансверсального обхода, который однозначно (с точностью до обращения) определяет граф и циклические порядки в вершинах. Таким образом, изучение комбинаторных свойств сборных графов сводится к анализу два-слов. Эта связь открывает богатые возможности для применения как теоретико-графовых, так и теоретико-языковых методов, а сама задача находит пересечения с теорией узлов, виртуальных узлов, родов поверхностей~\cite{Kauffman1999,Cairns1996}.

Существование простых сборных графов со сколь угодно большим сборным числом означает, что в общем случае нет никаких ограничений на количество генов, которые могут быть закодированы одной исходной последовательностью ДНК с заданной системой указателей (\cite{Angeleska2009}, Предложение 4.9). Нас интересуют два вопроса, связанные с экстремальными свойствами сборных чисел:

\begin{itemize}
    \item Задача о \emph{реализующем числе} $R_{\min}(k)$ — минимальном числе $4$-валентных вершин в простом сборном графе, имеющем сборное число $k$.
    \item Для заданного количества вершин $N$ найти максимально возможное значение $\An(\Gamma)$ и описать графы, на которых этот максимум достигается.
\end{itemize}

В работе~\cite{Angeleska2009} была предложена конструкция, дающая верхнюю оценку $R_{\min}(n)\le 3n-2$. Позднее в~\cite{Muche2012} было введено понятие графа \emph{петельной подстановки} (loop‑saturated graph) --- граф, в котором каждое исходное ребро заменено на петлю. Показано, что для таких графов выполняется $|\hGamma|=3\cdot\An(\hGamma)-2$. Совсем недавно в работе~\cite{MaksaevPreprint} с помощью анализа жадного алгоритма, строящего малое гамильтоново множество полигональных путей, была доказана нижняя оценка $R_{\min}(n)\ge 3n-2$: для простого сборного графа $\Gamma$ строится гамильтоново множество полигональных путей $\gamma$ такое, что $|\Gamma| \ge 3 \cdot |\gamma|  -2 \ge 3\cdot\An(\Gamma) - 2$. Гамильтоновы множества полигональных путей отождествляются с набором рёбер, которые в них входят: задача минимизации размера множества путей заменяется на задачу максимизации совокупного количества рёбер в путях. 
Аккуратный анализ алгоритма также позволяет установить, что максимум $\An(\Gamma)$ при заданном $|\Gamma|=N$ достигается в точности на графах петельных подстановок. Таким образом, обе упомянутые задачи полностью решены.

Поскольку петельные подстановки неизбежно содержат большое количество петель, естественно поставить вопрос о том, как изменится реализующее число при переходе к семейству \emph{простых сборных графов без петель}. Обозначим соответствующее число реализации через $R_{\min}^\circ(n)$. Методов, использованных в~\cite{MaksaevPreprint}, уже недостаточно, и для решения задачи требуется более тонкий инструментарий.

В настоящей работе мы развиваем подход, основанный на симметричной версии жадного алгоритма $\mathcal{G}_s$, и с его помощью получаем следующие основные результаты:

\begin{itemize}
	\item Для всякого простого сборного графа без петель справедливо неравенство
	\[
	|\Gamma| \ge 5\cdot \An(\Gamma)-3,
	\]
	откуда, аналогично общему случаю, следует нижняя оценка $R_{\min}^\circ(k)\ge5k-3$.
	\item Предъявлено полное семейство графов, на которых достигается равенство $|\Gamma|=5\cdot\An(\Gamma)-3$; оно состоит в точности из $(1212)$-подстановок произвольных (возможно, содержащих петли) сборных графов. Тем самым доказано, что $R_{\min}^\circ(n)=5n-3$.
	\item Дана комбинаторная характеризация экстремальных графов без петель: показано, что для любого графа, удовлетворяющего $\lfloor{\frac{|\Gamma| + 3}{5}}\rfloor = \An(\Gamma)$, соответствующее два-слово допускает разложение, описываемое регулярным выражением.
    \item Дана полная комбинаторная характеризация графов, удовлетворяющих $|\Gamma| = 5\cdot\An(\Gamma) - 2$ путём компьютерного перебора слов некоторого семейства автоматных языков.
\end{itemize}

Для получения этих результатов мы вводим понятие паттерна два-слова, сводящее исходную структуру к правильной скобочной последовательности, изучаем комбинаторику открывающих и закрывающих блоков и применяем технику пересечения недетерминированных конечных автоматов, ассоциированных с прямой и обращённой работой жадного алгоритма.

% Работа организована следующим образом. В~разделе~\ref{sec:term} собраны основные определения и обозначения. Раздел~\ref{sec:greedy} посвящён жадному алгоритму и его свойствам. В~разделе~\ref{sec:substitutions} исследуются $w$-подстановки и их влияние на сборное число. Основные теоремы о реализующем числе и характеризации экстремальных графов без петель доказаны в~разделе~\ref{sec:main}. Наконец, раздел~\ref{sec:conclusion} содержит обсуждение полученных результатов и возможные направления дальнейших исследований.

\section{Терминология и обозначения}\label{sec:term}

Мы придерживаемся обозначений и определений, введённых в~\cite{Angeleska2009,Muche2012,MaksaevPreprint}. В этой статье мы рассматриваем конечные связные неориентированные графы с вершинами степени 1 или 4. Петли и кратные рёбра допускаются. Для такого графа $\Gamma$ через $V(\Gamma)$ и $E(\Gamma)$ обозначаются множества вершин и рёбер между двумя 4-валентными вершинами графа $\Gamma$ соответственно. Повсюду в тексте под степенью вершины понимается число полурёбер, т.е. частей рёбер в локальной окрестности вершины, инцидентных этой вершине, при этом все кратные рёбра считаются по отдельности, а каждая петля считается дважды. Точнее говоря, если существует петля, инцидентная $v \in V(\Gamma)$, то она представляется двумя полурёбрами, добавляя тем самым два к степени $v$.

Для набора полурёбер $(e_1, e_2, e_3, e_4)$ мы определяем его \emph{циклический порядок} как следующее множество, состоящее из всех его циклических сдвигов и их обращений:
\begin{equation*}
\begin{aligned}
(e_1, e_2, e_3, e_4)_{\text{cyc}} = \{&(e_1, e_2, e_3, e_4),\,
(e_2, e_3, e_4, e_1),\,
(e_3, e_4, e_1, e_2),\,
(e_4, e_1, e_2, e_3),\\
&(e_4, e_3, e_2, e_1),\,
(e_3, e_2, e_1, e_4),\,
(e_2, e_1, e_4, e_3),\,
(e_1, e_4, e_3, e_2)\}.
\end{aligned}
\end{equation*}
Иными словами, мы считаем все элементы множества $(e_1, e_2, e_3, e_4)_{\text{cyc}}$ эквивалентными и обычно ссылаемся на один элемент этого порядка. Мы говорим, что вершина $v \in V(\Gamma)$ степени 4 является \emph{регулярной}, если циклический порядок полурёбер, инцидентных этой вершине, зафиксирован (вообще говоря, есть всего три способа сделать это). Повсюду в тексте мы изображаем графы как локально вложенные в поверхность в каждой вершине. Каждая вершина $v$ рассматривается как маленький диск, так что инцидентные рёбра прикреплены к границе диска. Циклический порядок полурёбер обычно задаётся так, как мы считываем их на диаграмме, следуя границе диска.

Предположим, что для вершины \( v \) степени 4 задан следующий циклический порядок полурёбер: \((e_1, e_2, e_3, e_4)\). Тогда мы говорим, что \( e_1 \) и \( e_2 \) являются \emph{соседними в \( v \)}, также как \( e_2 \) и \( e_3 \), \( e_3 \) и \( e_4 \), а также \( e_4 \) и \( e_1 \). Заметим, что полурёбра петли также могут быть соседними. Мы говорим, что \( e_1 \) и \( e_3 \), а также \( e_2 \) и \( e_4 \) являются \emph{трансверсальными в \( v \)}. Пара рёбер, инцидентных \( v \), является трансверсальной тогда и только тогда, когда рёбра не являются соседними.

\emph{Сборный граф} --- это конечный связный граф, в котором все вершины имеют степень 1 или 4, и все вершины степени 4 являются \emph{регулярными}. Вершину степени 1 будем называть \emph{концевой вершиной}.

\emph{Путь} в графе --- это чередующаяся последовательность вершин и рёбер, которая начинается и заканчивается вершиной, причём рёбра в этой последовательности инцидентны вершинам, стоящим рядом с ними. Если путь не содержит рёбер, то есть состоит из единственной вершины, он называется \emph{синглтоном}. Под длиной пути мы будем понимать число рёбер в нём.

\begin{definition}[\cite{Angeleska2009}, Определение 3.3] \emph{Трансверсальный путь} в сборном графе, или просто \emph{трансверсаль}, --- это путь \( v_0 e_1 v_1 e_2 \cdots e_n v_n \), где \( v_0, v_n \) --- концевые вершины, удовлетворяющий следующим условиям:

\begin{itemize}
    \item Все рёбра \( e_1, \ldots, e_n \) попарно различны.
    \item Каждое ребро \( e_i \) не является соседним с \( e_{i-1} \) в регулярной вершине \( v_{i-1} \) (\( i = 2, \ldots, n \)). Когда \( e_i \) является петлёй (\( v_{i-1} = v_i \)), оба полурёбра \( e_i \) помечаются как \( e_i^{(1)} \) и \( e_i^{(2)} \) таким образом, что \( e_{i-1} \) и \( e_i^{(1)} \) не являются соседними в \( v_i \), а также \( e_i^{(2)} \) и \( e_{i+1} \) не являются соседними.
\end{itemize}

\end{definition} 

Пример трансверсального пути можно видеть на Рисунке~2. Заметим, что в вершине \( v_1 \) трансверсальный путь сначала идёт по <<правому>> полурёбру, соответствующему петле \( e_1 \), а затем проходит через <<левое>> полурёбро петли \( e_1 \). В то же время в вершине \( v_3 \) путь сначала идёт по <<левому>> полурёбру петли \( e_4 \), а затем проходит через её <<правое>> полурёбро.

О сборном графе удобно думать как о \emph{семействе замкнутых гладких кривых} в трёхмерном пространстве с
$N$ различными точками, в которых две кривые (не обязательно различные) пересекаются по разным направлениям. Все такие точки пересечения соответствуют 4-валентным вершинам графа, а дуги кривых, соединяющие эти точки, --- рёбрам. В этой геометрической модели два полурёбра являются соседними, если их направления касательных в точке пересечения различны, а трансверсальный путь — это просто гладкая кривая.

Два трансверсальных пути эквивалентны, если они совпадают. \emph{Простым сборным графом} называется сборный граф, обладающий эйлеровым трансверсальным путём, т.е. трансверсальным путём, проходящим по каждому ребру ровно один раз. Заметим, что в простом сборном графе есть лишь два класса эквивалентности трансверсальных эйлеровых путей, которые отличаются обращением, см. \cite[Лемма 3.6]{Angeleska2009}. У простого сборного графа зафиксирована ориентация: одну из концевых вершин мы будем называть начальной, вторую --- финальной. Мы говорим, что два графа \emph{изоморфны как простые сборные графы}, если существует изоморфизм графов между ними, который переводит начальную концевую вершину в начальную концевую вершину другого графа и продолжается до биекции на полурёбрах, сохраняющей их циклический порядок во всех регулярных вершинах. Это означает, что простые сборные графы однозначно определяются трансверсальным эйлеровым путём, см. \cite[Определение 3.2 и Лемма 3.7]{Angeleska2009}.

Путь \( v_0e_1v_1 \cdots v_{\ell-1}e_\ell v_\ell \) в сборном графе называется \emph{полигональным} (см. \cite[Раздел 2]{6}), если все вершины \( v_0, \ldots, v_\ell \) являются попарно различными вершинами степени 4 графа \( \Gamma \), а \( e_i \) и \( e_{i+1} \) являются соседними в \( v_i \) для всех \( i = 1, \ldots, \ell - 1 \). Заметим, что цикл (в частности, петля) не является полигональным путём, поскольку все вершины в нём должны быть различны, в то время как изолированная вершина представляет собой полигональный путь. Примеры полигональных путей представлены на Рисунке 3.

Неформально, если рассматривать вершину как перекрёсток, мы можем либо повернуть налево или направо, образуя полигональный путь, либо проехать прямо в каждой вершине, что даёт нам трансверсальный путь.

Два пути \emph{не пересекаются по вершинам}, если они не имеют общих вершин. Попарно вершинно-непересекающееся множество полигональных путей в сборном графе \( \Gamma \) называется \emph{гамильтоновым}, если его объединение покрывает все вершины степени 4 графа \( \Gamma \) (см. Рисунок 3). Мы обозначаем через \( C(\Gamma) \) совокупность всех гамильтоновых множеств полигональных путей в \( \Gamma \).

Пусть $\gamma$ --- гамильтоново множество полигональных путей в простом сборном графе $\Gamma$. Через $E(\gamma)$ мы обозначаем множество рёбер, задействованных в полигональных путях. Заметим, что по $E(\gamma)$ можно однозначно восстановить $\gamma$. В частности, $E(\gamma) = \emptyset$ соответствует $\gamma$, покрывающему каждую 4-валентную вершину синглтоном.

\begin{lemma}[\cite{MaksaevPreprint}]\label{lemma:edge-setsize-rel}
    Пусть $\gamma = \{\gamma_1, \gamma_2, \dots,\gamma_t\}$ --- некоторое гамильтоново множество полигональных путей некоторого простого сборного графа $\Gamma$. Тогда $E(\gamma) = \bigcup_i\gamma_i$ --- множество рёбер путей из $\gamma$.
    \[
    |\Gamma| = |E(\gamma)| + |\gamma|
    \]
\end{lemma}
    
Действительно, полигональные пути формируют лес в графе, индуцированном $E(\gamma)$. Для леса это утверждение хорошо известно.

% Пусть \( \Sigma = \{a_1, a_2, \ldots\} \) --- (конечный) алфавит, а \( w = w_1w_2 \cdots w_{N-1}w_N \) --- слово длины \( N \) в алфавите \( \Sigma \). Длина \( w \) обозначается через \( |w| \), так что \( |w| = N \). Мы полагаем \( w[i] = w_i \), т.е. \( i \)-й символ слова \( w \). Через \( w[i : j] \) (\( i \leq j \)) мы обозначаем подслово \( w_i w_{i+1} \cdots w_{j-1} w_j \), при этом \( |w[i : j]| = j - i + 1 \). Пусть \( a \in \Sigma \) встречается ровно дважды в фиксированном слове \( w \) на позициях \( i \) и \( j \), т.е. \( w[i] = w[j] = a \), \( i < j \). Мы определяем две функции \( \sigma_1 \) и \( \sigma_2 \), которые дают первое и второе вхождение \( a \) в \( w \) соответственно, т.е. \( \sigma_1(a) = i \) и \( \sigma_2(a) = j \). В дальнейшем использование \( \sigma_1(a) \) или \( \sigma_2(a) \) предполагает, что слово \( w \) ясно из контекста и буква \( a \) появляется в \( w \) ровно дважды.

% \emph{сборным словом}, или \emph{словом с двукратным вхождением} (DOW), в алфавите \( \Sigma \) называется непустое слово, в котором каждый символ алфавита встречается либо ровно дважды, либо не встречается вовсе. Заметим, что это комбинаторное определение используется для гауссовых слов в

Пусть $\Sigma=\{a_1,a_2,\ldots\}$ — конечный алфавит. Слово $w=w_1w_2\ldots w_{2n}$ из $2n$ букв в алфавите $\Sigma$ называется \emph{два-словом} или сборным словом (double occurrence word, DOW), если каждая буква из $\Sigma$ встречается в нём ровно два раза. Число $n=|w|$ называют длиной два-слова.

Отметим, что это комбинаторное определение используется для гауссовых слов в теории плоских кривых. Два слова называются \textit{эквивалентными}, если одно из другого можно получить переименованием символов алфавита. Существует взаимно-однозначное соответствие между классами изоморфизма простых сборных графов и классами эквивалентности два-слов, как доказано в \cite[Лемма 3.8]{Angeleska2009}. Ниже мы кратко приводим эту конструкцию для полноты и дальнейшего использования.

Перечисление вершин степени 4, посещаемых трансверсальным путём в простом сборном графе, даёт два-слово над алфавитом \(V(\Gamma)\), поскольку любой трансверсальный путь должен посетить каждую вершину степени 4 дважды. Следовательно, сборные слова эквивалентны тогда и только тогда, когда они соответствуют изоморфным сборным графам, как утверждается.

Для простого сборного графа \(\Gamma\) обозначим представителя соответствующего класса сборных слов через \(w_\Gamma\). Обратно, \emph{простой сборный граф} $\Gamma_w$, ассоциированный с два-словом $w$, строится следующим образом.
\begin{itemize}
    \item Множество вершин $V(\Gamma_w)$ состоит из всех различных букв слова $w$; все эти вершины считаются $4$-валентными.
    \item Рёбра графа отвечают последовательным парам букв слова: для каждого $i=1,\dots,2n-1$ вводится ребро $e_i=(w_i,w_{i+1})$.
    \item Циклический порядок полурёбер в каждой вершине задаётся порядком обхода вдоль слова $w$ (эйлеров трансверсальный путь). Иными словами, в локальной окрестности вершины полурёбра упорядочены так, как они появляются при движении от первого вхождения буквы ко второму и далее к концу слова.
    \item В слово дополнительно включают два фиктивных концевых символа, дающих начало и конец трансверсального пути; соответствующие им вершины имеют степень $1$ и называются \emph{концевыми}.
\end{itemize}
Граф $\Gamma_w$ является связным, все его $4$-валентные вершины жёсткие, а концевые вершины имеют степень $1$.

\emph{Реализующим числом} $R_{\min}(k)$ называется минимальное количество $4$-валентных вершин в простом сборном графе, сборное число которого равно $k$:
\[
R_{\min}(k)=\min\{\,|\Gamma|:\An(\Gamma)=k\,\}.
\]
Аналогично, $R_{\min}^\circ(k)$ обозначает реализующее число в классе простых сборных графов без петель.

Пусть $w$ и $u$ — два два-слова без общих букв. Их \emph{композицией} $w\circ u$ называется конкатенация $wu$. Композиция простых сборных графов $\Gamma_w$ и $\Gamma_u$ обозначается как $\Gamma_w \circ \Gamma_u$ и соответствует два-слову композиции два-слов $w\circ u$.

\begin{definition}[{\cite[Definition 2.3.1]{Muche2012}}]
    Сборный граф $\Gamma$ называется \emph{аддитивным слева}, если для любого сборного графа $\Gamma'$ выполнено
    \[
        \An(\Gamma \circ \Gamma') = \An(\Gamma) + \An(\Gamma').
    \]
    Аналогично, $\Gamma$ называется \emph{аддитивным справа}, если для любого $\Gamma'$ выполнено
    \[
        \An(\Gamma' \circ \Gamma) = \An(\Gamma') + \An(\Gamma).
    \]
    Граф, являющийся одновременно аддитивным слева и справа, называется \emph{аддитивным}.

    Сборный граф $\Gamma$ называется \emph{двусторонне аддитивным}, если для любых сборных графов $\Gamma_1,\Gamma_2$ справедливо
    \[
        \An(\Gamma_1 \circ \Gamma \circ \Gamma_2) =
        \An(\Gamma_1) + \An(\Gamma) + \An(\Gamma_2).
    \]
\end{definition}

\begin{lemma}[{\cite[Лемма 2.3.6]{Muche2012}}]\label{lem:additivity-criterion}
    Пусть $\Gamma$ — сборный граф, а $\Gamma_{(1)}$ — граф петли, соответствующий два-слову $(11)$. Тогда
    \begin{enumerate}
        \item $\Gamma$ аддитивен справа $\iff$ $\An(\Gamma_{(1)} \circ \Gamma) = \An(\Gamma) + 1$;
        \item $\Gamma$ аддитивен слева $\iff$ $\An(\Gamma \circ \Gamma_{(1)}) = \An(\Gamma) + 1$;
        \item $\Gamma$ двусторонне аддитивен $\iff$ $\An(\Gamma_{(1)} \circ \Gamma \circ \Gamma_{(1)}) = \An(\Gamma) + 2$.
    \end{enumerate}
\end{lemma}

Всякий двусторонне аддитивный граф является аддитивным, но не наоборот. Граф аддитивен, но не двусторонне аддитивен, когда справедливо:

\begin{enumerate}
    \item $\An(\Gamma_{(1)} \circ \Gamma) = \An(\Gamma) + 1$;
    \item $\An(\Gamma \circ \Gamma_{(1)}) = \An(\Gamma) + 1$;
    \item $\An(\Gamma_{(1)} \circ \Gamma \circ \Gamma_{(1)}) = \An(\Gamma) + 1$.
\end{enumerate}

Мы называем граф \emph{неаддитивным}, если хотя бы одно условие аддитивности для него нарушается: $\An(\Gamma_{(1)} \circ \Gamma) = \An(\Gamma)$ либо $\An(\Gamma \circ \Gamma_{(1)}) = \An(\Gamma)$.

\begin{definition}
	\emph{Петельной подстановкой} $\hGamma$ простого сборного графа $\Gamma$ называется граф, полученный заменой каждого ребра графа $\Gamma$ петлёй (что соответствует подстановке два-слова $11$ на место каждого ребра). В более общем случае, для произвольного непустого два-слова $w$ определяется \emph{$w$-подстановка} $\mathcal{S}_w(\Gamma)$ — граф, получающийся заменой каждого ребра графа $\Gamma$ копией графа $\Gamma_w$.
\end{definition}

Наряду с графом $\Gamma$ мы часто рассматриваем обращённый граф $\Gamma^R$, соответствующий два-слову, прочитанному в обратном порядке. Сборное число при обращении не меняется: $\An(\Gamma)=\An(\Gamma^R)$.

\section{Подстановки}

Для изучения простых сборных графов без петель нам понадобится обобщить определение петельной подстановки.

\begin{definition}
     Сборный граф $\mathcal{S}_w(\Gamma)$, полученный из данного сборного графа $\Gamma$ и два-слова $w$ путём замены всех рёбер графа $\Gamma$ на копии графа $\Gamma_{w}$, будем называть \textit{сборным графом, полученным из $\Gamma$ путём $w$-подстановки}.

     Сборный граф $\mathcal{S}^{\circ}_w(\Gamma)$, полученный из данного сборного графа $\Gamma$ и два-слова $w$ путём замены всех рёбер графа $\Gamma$ на копии графа $\Gamma_{w}$, кроме рёбер, инцидентных концевым вершинам, будем называть \textit{сборным графом, полученным из $\Gamma$ путём внутренней $w$-подстановки}.
\end{definition}

Таким образом, определения петельной подстановки $\hGamma$ и внутренней петельной подстановки $\hGamma^\circ$ из \cite{Muche2012} совпадают с $\mathcal{S}_{(11)}(\Gamma)$ и $\mathcal{S}^{\circ}_{(11)}(\Gamma)$.

Аналогично \cite{MaksaevPreprint}, мы отождествляем гамильтоново множество полигональных путей $\gamma$ с множеством рёбер, которые в него входят $E(\gamma)$. Часто мы не будем различать задачу минимизации числа полигональных путей и задачу максимизации числа рёбер в них, так как, уже было сказано ранее, они связаны линейным уравнением.
\[
    |\Gamma| = |\gamma| + E(\gamma)
\]
Чтобы проверить, что какое-то подмножество рёбер $E' \subset E(\Gamma)$ задаёт валидное гамильтоново множество полигональных путей, нужно проверить два условия.
\begin{itemize}
    \item \emph{Полигональность} --- никакие два последовательных ребра трансверсального пути одновременно не находятся в $E'$. 
    \item \emph{Отсутствие циклов} --- в ограничении $\Gamma$ на $E'$ нет циклов. 
\end{itemize}

Подобно множеству задач из теории графов, мы будем называть процесс добавления ребра в $E(\gamma)$ \emph{покраской} ребра. 

Пусть дано слово $w$ длины $m$ и простой сборный граф $\Gamma$ с $n$ 4-валентными вершинами, которому соответствует два-слово $g$. Тогда графу $\check{\Gamma} = \mathcal{S}_w(\Gamma)$ соответствует слово:
\[
w^1,g_1,w^2,g_2,\dots,w^{2n},g_{2n},w^{2n + 1}
\]
где $\{w^i\}_{i = 1}^{2n + 1}$ --- копии слова $w$. 

Введём обозначения для следующих подмножеств рёбер.

\begin{enumerate}
    \item $E_{w^i}$ --- внутренние рёбра $w^i$
    \[
        E_{w^i}= \{e \in E(\check{\Gamma}) : V(e) \subset V(w^i)\}
    \]
    \item $E_{g}$ --- рёбра, связывающие копии $w$ с исходными вершинами $g$
    \[
        E_{g} = \{e \in E(\check{\Gamma}): \exists g_i \in V(e)\}. 
    \]
\end{enumerate}

Заметим, что каждое ребро $\check{\Gamma}$ попало ровно в одно подмножество:  $E(\check{\Gamma}) = (\bigsqcup_i^{2n + 1} E_{w^i}) \sqcup E_g$.

\begin{theorem}\label{theor:an-saturation} 
    Для любого простого сборного графа $\Gamma$ и любого непустого два-слова $w$ верны равенства:
    
    \begin{itemize}

        \item $\An(\mathcal{S}_w(\Gamma)) \leq (2|\Gamma| + 1)\An(w) + |\Gamma|$
    
        \item $\An(\mathcal{S}_w(\Gamma)) \geq (2|\Gamma| + 1)\An(w) - |\Gamma| $.
    \end{itemize}
\end{theorem}


\begin{proof}

    Докажем каждое неравенство по отдельности.

    \begin{enumerate} 
        \item Пусть $\gamma_w$ будет некоторым минимальным гамильтоновым множеством полигональных путей для $w$. Множество $\gamma_w$ индуцирует покраску рёбер всякого $w^i$. В совокупности это задаёт валидное гамильтоново множество полигональных путей $\gamma$ для $\check{\Gamma}$. 
        \[ 
        E(\gamma) \cap E_{w^i} \sim E(\gamma_w) \quad \forall i 
        \]
        Полигональность выполняется по построению. Граф $\check{\Gamma}$, ограниченный на $E(\gamma)$, разбивается на $2n + 1$ подграфов, изоморфных графу $\Gamma_w$, ограниченному на $E(\gamma_w)$. Ни в одном из них циклов нет, а значит $\gamma$ --- корректно. Итого $\An(\check{\Gamma}) \leq |\gamma| = (2|\Gamma| + 1) \cdot \An(w) + |\Gamma|$

        \item В силу полигональности не больше $2n$ рёбер из $E_g$ могут единовременно быть покрашены. С другой стороны, покраска любого $w^i$ должна задавать корректную покраску $w$. 

        \[
        |E(\gamma)| = \sum_i|E(\gamma)\cap E_{w^i}| + |E(\gamma) \cap E_g| \leq (2n + 1)(m - \An(w)) + 2n
        \]
        
        Отсюда получаем $\An(\check{\Gamma}) = |\check{\Gamma}| - |E(\gamma)| \geq (2n + 1)\cdot \An(w) - n$

    \end{enumerate}
\end{proof}

Сборное число $w$-подстановки тесно связано с понятиями аддитивности и двусторонней аддитивности. 

\begin{lemma}\label{lemma:left-additive}
    Простой сборный граф $\Gamma$ является аддитивным слева тогда и только тогда, когда не существует минимального гамильтонова множества полигональных путей такого, что заключительное ребро его трансверсального пути не закрашено.
\end{lemma}

\begin{proof}
    Как уже было сказано ранее, простой сборный граф $\Gamma$ является аддитивным слева, если $\An(\Gamma\circ \Gamma_{(1)}) = \An(\Gamma) + 1$. Закрашивать петлю $\Gamma_{(1)}$ нельзя, так как это создаст цикл. Пусть $e$ --- ребро, соединяющее $\Gamma$ с петлей $\Gamma_{(1)}$. Пусть $\gamma$ --- минимальное гамильтоново множество полигональных путей для графа $\Gamma\circ \Gamma_{(1)}$. Разберём два случая.

    \begin{itemize}
        \item Ребро $e$ не закрашено. В таком случае $E(\gamma)$ задаёт корректное гамильтоново множество полигональных путей для $\Gamma$. Следовательно, $\An(\Gamma\circ\Gamma_{(1)}) = |\Gamma| - |E(\gamma)| + 1 = \An(\Gamma) + 1$.
        \item Ребро $e$ закрашено. В таком случае $E(\gamma) \setminus \{e\}$ задаёт корректное гамильтоново множество полигональных путей для $\Gamma$ (назовём его $\gamma'$). Причём, в силу полигональности, заключительное ребро трансверсального пути $\Gamma$ у $\gamma'$ не закрашено. Нам интересен лишь случай $|E(\gamma)\setminus \{e\}| = |\Gamma| - \An(\Gamma)$, поскольку в противном случае мы получаем множество полигональных путей не лучше, чем если бы не закрашивали $e$. Если $E(\gamma) - 1 = |\Gamma| - \An(\Gamma)$, мы получаем $\An(\Gamma\circ\Gamma_{(1)}) = |\Gamma| - |E(\gamma)| + 1 = \An(\Gamma)$.
    \end{itemize}

    Если у исходного графа $\Gamma$ существует минимальное гамильтоново множество полигональных путей $\gamma$, которое не закрашивает последнее ребро трансверсального пути, то $E(\gamma) \cup \{e\}$ задаёт корректное гамильтоново множество полигональных путей для $\Gamma\circ \Gamma(1)$ размера $\An(\Gamma)$. Из рассуждения выше понятно, что меньше множества найти нельзя. В противном случае, как и было заявлено, сборное число $\An(\Gamma\circ \Gamma_{(1)}) = \An(\Gamma) + 1$.
\end{proof}

\begin{lemma}\label{lemma:middle-additive}
    Простой аддитивный сборный граф $\Gamma$ является двусторонне аддитивным тогда и только тогда, когда не существует гамильтонова множества полигональных путей размера $\An(\Gamma) + 1$ такого, что первое и заключительное ребро его трансверсального пути одновременно не закрашены.
\end{lemma}

\begin{proof}
    Доказательство полностью аналогично рассуждению из леммы \ref{lemma:left-additive}.
\end{proof}

\begin{conjecture}\label{conj:an-additive-saturation}
    Пусть дан простой сборный граф $\Gamma$ и непустое неаддитивное два-слово $w$. Тогда сборное число $w$-подстановки имеет следующий вид.
    \[
        \An(\mathcal{S}_w(\Gamma)) = (2|\Gamma| + 1)\cdot\An(w) - |\Gamma|
    \]

\end{conjecture}

\begin{proof}
    Формулу достаточно доказать либо для $w$, либо для $w^R$, поскольку сборное число не меняется при обращении слова. Заметим, что аддитивность справа у $w$ эквивалентна аддитивности слева у $w^R$. Следовательно, неаддитивность $w$ означает, что либо $w$, либо $w^R$ не является аддитивным слева. Без ограничения общности будем считать, что аддитивность слева нарушена именно у $w$.
 
    У два-слова $w$ существует минимальное гамильтоново множество полигональных путей $\gamma_w$ с незакрашенным последним ребром трансверсального пути \ref{lemma:left-additive}.
    Построим множество полигональных путей $\gamma$ для $\check{\Gamma}$ следующим образом.
    \begin{itemize}
        \item Покраска $w^i$ индуцирована множеством $E(\gamma_w)$.
        \[
            E(\gamma) \cap E_{w^i} \sim E(\gamma_w)
        \]
        \item Пусть рёбра $e_l^i, e_r^i$ соединяют вершину $g_i$ с $w^i$ и $w^{i + 1}$ соответственно. Закрасим все рёбра $\{e_l^i\}$.
    \end{itemize}

    Чтобы убедиться в корректности $\gamma$, нужно проверить два свойства.

    \begin{itemize}
        \item Полигональность выполняется по построению. Рёбра $\{e_l^i\}$ ни с чем не конфликтуют благодаря выбору $\gamma_w$. 
        \item Отсутствие циклов. Предположим противное: в $E(\gamma)$ есть цикл положительной длины. С одной стороны, цикл не может целиком содержаться ни в одной компоненте $w^i$ по построению. С другой стороны, цикл не может проходить положительное число раз через разрез, определяемый вершинами какой-либо компоненты $V(w^i)$, поскольку любой цикл обязан проходить через разрез чётное число раз, а множество $E(\gamma)$ пересекается с таким разрезом лишь по ребру $e^i_l$. Мы пришли к противоречию.
    \end{itemize} 
     Мы закрасили $2n$ рёбер из $E_g$ и $m - \An(w)$ рёбер из каждого $E_{w^i}$. Из доказательства второго пункта теоремы \ref{theor:an-saturation} понятно, что это максимальное количество рёбер, которое может быть закрашено в каждой из компонент. Следовательно, это минимальное гамильтоново множество полигональных путей.
    \[
    \An(\mathcal{S}_w(\Gamma)) = (2|\Gamma| + 1)\cdot\An(w) - |\Gamma|
    \]

    % $(\Longrightarrow)$ Пусть $w$ --- аддитивно. Это значит, что у любого минимального гамильтонового множества полигональных путей $\gamma_w$ для $w$ крайние рёбра $(w^i_1, w^i_2), (w^i_{2m - 1}, w^i_{2m}) \in E(\gamma_w)$. 
    
    % Пусть для $\hGamma$ существует множество полигональных путей $\gamma$ размера $(2n + 1)\cdot \An(w) - n$. Это возможно только если неравенства на компоненты $E(\gamma) \cap E_{w^i} \text{ и } E(\gamma) \cap E_g$, упомянутые в доказательстве второго пункта \ref{theor:an-saturation}, превратились в равенства. 
    % \[
    % (|E(\gamma) \cap E_g| = 2n )\wedge (|E(\gamma) \cap E_{w^i}| = m - \An(w) \quad \forall i)
    % \]

    % Но это значит, что каждая компонента $E(\gamma) \cap E_{w^i} \sim E(\gamma_w) \quad \forall i$ для некоторого минимального множества полигональных путей $\gamma_w$. Следовательно, полигональность не будет нарушена только если $E(\gamma) \cap E_g = \emptyset$. Противоречие.
\end{proof}

\begin{corollary}\label{cor:pretzelsat5k-3}
    Пусть $\Gamma$ простой сборный граф размера $k-1$.
    \begin{itemize}
        \item $(1212)$-подстановка графа $\Gamma$ имеет размер $5k - 3$ и сборное число $k$.
        \item $(11)$-подстановка (петельная подстановка) графа $\Gamma$  имеет размер $3k - 2$ и сборное число $k$.
    \end{itemize}
\end{corollary}

\begin{proof}

Действительно, два-слова $(1212)$ и $(11)$ являются неаддитивными:
\[
    \An(\Gamma_{(11)} \circ \Gamma_{(1212)}) = \An(\Gamma_{(1212)} \circ \Gamma_{(11)}) =\An(\Gamma_{(1212)}) = 1.
\]

\[
    \An(\Gamma_{(11)} \circ \Gamma_{(11)}) = \An(\Gamma_{(11)}) = 1.
\] 
А значит, из теоремы~\ref{conj:an-additive-saturation} следует, что $\An(\mathcal{S}_{(1212)}(\Gamma)) = \An(\mathcal{S}_{(11)}(\Gamma)) = (2\cdot |\Gamma| + 1) -|\Gamma| = |\Gamma| + 1 = k$. Размеры подстановок имеют следующий вид.

\[
    |\mathcal{S}_{(11)}(\Gamma)| = (2\cdot |\Gamma| + 1) \cdot 1 + |\Gamma| = 3\cdot |\Gamma| + 1 = 3k - 2.
\]

\[
    |\mathcal{S}_{(1212)}(\Gamma)| = (2\cdot |\Gamma| + 1) \cdot 2 + |\Gamma| = 5\cdot |\Gamma| + 2 = 5k - 3.
\]

\end{proof}

\begin{definition}
    Пусть $\Gamma$ --- простой сборный граф. Определим функцию $\varphi_{i,j}(\Gamma),\; i,j \in \{0,1\}$ как максимальное количество закрашенных рёбер среди всех гамильтоновых множеств полигональных путей для $\Gamma$, таких что выполняются следующие условия.

    \begin{itemize}
        \item первое ребро трансверсального пути $\Gamma$ \emph{не закрашено}, если $i = 1$.
        \item последнее ребро трансверсального пути $\Gamma$ \emph{не закрашено}, если $j = 1$.
    \end{itemize}
\end{definition}

\begin{lemma}
    Пусть $\Gamma$ --- аддитивный простой сборный граф.
    \begin{itemize}
        \item $\varphi_{0,0}(\Gamma) = |\Gamma| - \An(\Gamma)$
        \item $\varphi_{0,1}(\Gamma) = \varphi_{1,0}(\Gamma) = |\Gamma| - \An(\Gamma) - 1$
        \item $\varphi_{1,1}(\Gamma) = |\Gamma| - \An(\Gamma) - 2$ если $\Gamma$ --- двусторонне аддитивный и $|\Gamma| - \An(\Gamma) - 1$ иначе.
    \end{itemize}
\end{lemma}

\begin{proof}
    Докажем каждый пункт по отдельности.
    \begin{itemize}
        \item Никаких ограничений нет. Следует напрямую из \ref{lemma:edge-setsize-rel}.
        \item В силу симметрии, достаточно доказать утверждение лишь для $\varphi_{0,1}$. Оценку снизу $\varphi_{1,0}(\Gamma) \geq |\Gamma| - \An(\Gamma) - 1$ можно получить, взяв любое минимальное гамильтоново множество полигональных путей для $\Gamma$ и удалив из него последнее ребро трансверсального пути (если оно там есть). В худшем случае число рёбер уменьшится на единицу. Чтобы получить оценку сверху, предположим обратное: существует гамильтоново множество полигональных путей $\gamma$ для $\Gamma$, у которого не закрашено последнее ребро. Из леммы \ref{lemma:left-additive} следует, что $\Gamma$ не является аддитивным слева. Противоречие.
        \item Оценку снизу $\varphi_{1,1}(\Gamma) \geq |\Gamma| - \An(\Gamma) - 2$ можно получить способом, аналогичным предыдущему пункту. Эта функция накладывает больше ограничений, чем все предыдущие, а значит $\varphi_{1,1}(\Gamma) \leq \varphi_{0,1}(\Gamma) = |\Gamma| - \An(\Gamma) - 1$. Из леммы \ref{lemma:middle-additive} непосредственно следует, что $\varphi_{1,1}(\Gamma) = |\Gamma| - \An(\Gamma) - 2$ исключительно для двусторонне аддитивных простых сборных графов. 
    \end{itemize}
\end{proof}

\begin{conjecture}
    Для любого простого сборного графа $\Gamma$ и двусторонне аддитивного два-слова $w$:
    \[\An(\mathcal{S}_w(\Gamma)) = (2|\Gamma| + 1) \cdot \An(w) + |\Gamma|\].
\end{conjecture}
\begin{proof}
    Рассмотрим рёбра из следующего отрезка трансверсального пути $\Gamma$.
    \[
        g_i, w^i, g_{i + 1}
    \]
    Для полноты изложения будем считать, что $g_0$ и $g_{2n + 1}$ равны каким-то фиктивным символам. Помимо рёбер из $E_{w^i}$ в таком блоке участвуют рёбра $e_i^r$ и $e_{i + 1}^l$. Максимальное количество рёбер блока, которое можно одновременно покрасить, не больше $\varphi_{0,0}(w), \varphi_{1,0}(w) + 1, \varphi_{0,1}(w) + 1$ либо $\varphi_{1,1}(w) + 2$ в зависимости от того, покрашены рёбра $e_i^r, e_{i + 1}^l$ или нет. Заметим, что для двусторонне аддитивных слов все эти четыре функции равны $|w| - \An(w)$. Столько рёбер всегда можно покрасить, не задействуя рёбра $e_i^r$ и $e_{i + 1}^l$. Следовательно, для таких подстановок гамильтоново множество полигональных путей, описанное в первом пункте \ref{theor:an-saturation}, является минимальным.
\end{proof}

\begin{conjecture}
    Для любого простого сборного графа $\Gamma$ и аддитивного, но не двусторонне аддитивного два-слова $w$:
    \begin{itemize}
        \item $\An(\mathcal{S}_w(\Gamma)) = (2|\Gamma| + 1) \cdot \An(w) + \An(\Gamma)$ если $\An(\mathcal{S}_w^\circ(\Gamma_{(11)})) = \An(w) + 1$;
        \item $\An(\mathcal{S}_w(\Gamma)) = (2|\Gamma| + 1) \cdot \An(w)$ если $\An(\mathcal{S}_w^\circ(\Gamma_{(11)})) = \An(w)$.
    \end{itemize}
\end{conjecture}

\begin{proof}
    Пусть $\gamma$ --- минимальное гамильтоново множество полигональных путей для $\mathcal{S}_w(\Gamma)$. Аналогично предыдущему рассуждению, разобьём трансверсальный путь на следующие блоки.
    \[
        g_i, w^i, g_{i + 1}
    \]
    Максимальное количество рёбер блока, которое можно одновременно покрасить, не больше $\varphi_{0,0}(w), \varphi_{1,0}(w) + 1, \varphi_{0,1}(w) + 1$ либо $\varphi_{1,1}(w) + 2$ в зависимости от того, покрашены рёбра $e_i^r, e_{i + 1}^l$ или нет. Для аддитивных, но не двусторонне аддитивных слов $w$ справедливо: $\varphi_{0,0}(w) = \varphi_{1,0}(w) + 1 = \varphi_{0,1}(w) + 1 = |w| - \An(w)$, в то время как $\varphi_{1,1}(w) = |w| - \An(w) + 1$. Как говорилось ранее, $|w| - \An(w)$ является базовым минимумом. Его можно достичь, не закрашивая рёбра $e_i^r$ и $e_{i + 1}^l$. Без ограничения общности будем считать, что $e_i^r$ и $e_{i + 1}^l$ либо одновременно закрашены, либо одновременно не закрашены. Пусть $S$ --- множество блоков, у которых закрашены одновременно $e_i^r$ и $e_{i + 1}^l$. Каждому блоку из $S$ биективно соответствуют два закрашенных ребра из $E_g$. Отсюда получаем оценку $2\cdot|S| \leq |E_g\cap E(\gamma)| \leq 2n$. Это даёт оценку снизу на сборное число.
    \[
        \An(\mathcal{S}_w(\Gamma)) \geq (2|\Gamma| + 1) \cdot \An(w)
    \]
    Перейдём к разбору случаев.
    \begin{itemize}
        \item Пусть существует гамильтоново множество полигональных путей $\gamma_w$ для $w$ размера $|w| - \An(w) - 1$, такое что первое и последнее ребро трансверсального пути не закрашены и первая и последняя вершина трансверсального пути не соединены одним полигональным путём. Это означает, что для графа внутренней петельной подстановки $\mathcal{S}_w^\circ(\Gamma_{(11)})$ можно закрасить $\varphi_{11}(w)$ рёбер подграфа, изоморфного $w$, а также первое и последнее ребро трансверсального пути. Эта конструкция корректна и оптимальна. Следовательно, $\An(\mathcal{S}_w^\circ(\Gamma_{(11)})) = \An(w)$.

        Построим гамильтоново множество полигональных путей $\gamma$ для $\mathcal{S}_w(\Gamma)$ следующим образом.
        \begin{itemize}
            \item Покраска $w^i$ индуцирована множеством $E(\gamma_w)$.
            \item Красим $e_r^i$ и $e_l^{i + 1}$ для всех нечётных $i$.
        \end{itemize}
        Для каждого нечётного $i$ блок находится в $S$. Такое $\gamma$ имеет минимальный возможный размер: $\An(\mathcal{S}_w(\Gamma)) = (2|\Gamma| + 1) \cdot \An(w)$.
        \item Пусть для любого гамильтонова множества полигональных путей $\gamma_w$ для $w$ размера $|w| - \An(w) - 1$, такого что первое и последнее ребро трансверсального пути не закрашены, существует полигональный путь, соединяющий первую и последнюю вершину трансверсального пути. Пусть $\gamma_w$ --- какое-то из них. Конструкция для $\mathcal{S}_w^\circ(\Gamma_{(11)})$ из предыдущего пункта уже не подходит, поскольку в этом случае в ней будет цикл. Поэтому сборное число внутренней подстановки на единицу больше: $\An(\mathcal{S}_w^\circ(\Gamma_{(11)})) = \An(w) + 1$.

        Каждый блок из $S$ соединяет полигональным путём вершины $g_i$ и $g_{i + 1}$. У сборного числа подстановки начинает зависеть от сборного числа простого сборного графа $\Gamma$:  $|S| \leq |\Gamma| - \An(\Gamma)$ и $\An(\mathcal{S}_w(\Gamma)) = (2|\Gamma| + 1) \cdot \An(w) + \An(\Gamma)$.
    \end{itemize}
\end{proof}


\section{Жадный алгоритм}

Дальше речь пойдёт о жадном алгоритме $\mathcal{G}$, который по простому сборному графу $\Gamma$ строит достаточно малое гамильтоново множество полигональных путей $\gamma = \G(\Gamma)$. В случае графа без петель мы получим оценку сверху $|\Gamma| \geq 5\cdot |\G(\Gamma)| - 3 \geq 5 \cdot \An(\Gamma) - 3$.

\begin{definition}
    Пусть $\Gamma$ --- простой сборный граф, где $e_1,e_2,\ldots,e_{2n - 1}$ его рёбра в порядке трансверсального обхода. Обозначим через $\G$ итеративный алгоритм, который по $\Gamma$ строит гамильтоново множество полигональных путей следующим образом:
   
    \begin{itemize}
        \item Начинаем с тривиального множества $\gamma$, $E(\gamma) = \emptyset$.
        \item На $i$-м шаге обрабатываем ребро $e_i$. Если после добавления ребра $e_i$ в $E(\gamma)$ $\gamma$ остаётся гамильтоновым множеством полигональных путей, добавляем его. 
    \end{itemize}
\end{definition}


\begin{lemma}
    Пусть $\Gamma$ --- простой сборный граф, где $e_1,e_2,\ldots,e_{2n - 1}$ его рёбра в порядке трансверсального обхода, а $\gamma$ --- некоторое гамильтоново множество полигональных путей. Пусть существуют рёбра $e_i, e_{i + 1} \not \in E(\gamma)$ такие, что:

    \begin{itemize}
        \item $i = 1$ либо $e_{i - 1} \not \in E(\gamma)$
        \item $i = 2n - 2$ либо $e_{i + 2} \not \in E(\gamma)$
    \end{itemize}    

    Тогда существует $e \in \{e_i, e_{i + 1}\}$ такое, что добавление $e$ к $E(\gamma)$ оставляет $\gamma$ гамильтоновым множеством полигональных путей.  
\end{lemma}

\begin{proof}
    preprint
\end{proof}

\begin{corollary}
    Пусть $\Gamma$ --- простой сборный граф, где $e_1,e_2,\ldots,e_{2n - 1}$ его рёбра в порядке трансверсального обхода, $\gamma$ --- множество полигональных путей, полученное алгоритмом $\G$. $E(\gamma) = \{e_{i_1}, e_{i_2},\ldots,e_{i_k}\}$ где $i_1 < i_2 <\ldots<i_k$. Тогда:
    \[i_{j + 1} = i_j + 2 \text{ либо } i_{j + 1} = i_j + 3 \quad \forall j < k\]
\end{corollary}

Теперь мы полностью перейдём к графам без петель. Далее будем считать, что все графы не имеют петель, если не сказано иное.  Определим два семейства блоков: $A = \{(e_{i_j}, e_{i_j + 1}) : i_{j} + 2 = i_{j + 1}  \vee i_j + 2 = 2n\}$ и $B = \{(e_{i_j}, e_{i_j + 1}, e_{i_j + 2}) : i_{j} + 3 = i_{j + 1} \vee i_j + 3 = 2n\}$

\begin{conjecture}\label{conj:no-loops-G-properties}
    Для графов без петель алгоритм $\G$ имеет следующие свойства:
    \begin{itemize}
        \item $i_1 = 1$
        \item $(\bigsqcup_{a \in A} a) \sqcup (\bigsqcup_{b \in B} b) = E(\Gamma)$ если $e_{2n - 1} \not \in E(\gamma)$
        \item $(\bigsqcup_{a \in A} a) \sqcup (\bigsqcup_{b \in B} b) = E(\Gamma) \setminus \{e_{2n - 1}\}$ если $e_{2n - 1} \in E(\gamma)$
    \end{itemize}
\end{conjecture}


Попробуем выразить имеющиеся параметры через $\alpha  = |A|$ и $\beta = |B|$.

\begin{itemize}
    \item Случай $e_{2n - 1} \not\in E(\gamma)$:
    \[
    \left\{
    \begin{aligned}
        & 2|\Gamma| - 1 = 2\alpha + 3\beta, \\
        & |E(\gamma)| = \alpha + \beta, \\
        & |\Gamma| = |\gamma| + |E(\gamma)|
    \end{aligned} 
    \right. 
    \implies \ 2\cdot|\gamma| = \beta + 1
    \]
    \item Случай $e_{2n - 1} \in E(\gamma)$:
    \[
    \left\{
    \begin{aligned}
        & 2|\Gamma| - 1 = 2\alpha + 3\beta + 1, \\
        & |E(\gamma)| = \alpha + \beta + 1, \\
        & |\Gamma| = |\gamma| + |E(\gamma)|
    \end{aligned} 
    \right.
    \implies \ 2\cdot|\gamma| = \beta
    \]
\end{itemize}

Видно, что из ограничения сверху на $\beta$ мы получаем ограничение сверху на $|\gamma|$. Дальше мы покажем, что $\beta \leq \alpha + 1$.


\begin{definition}
    Пусть буква $x$ появляется в два-слове $w = w_1,\ldots,w_{2n}$ на позициях $j_1$ и $j_2$, где $j_1 < j_2$. Будем говорить, что $w_{j_1}$ является первым появлением $x$, а $w_{j_2}$ вторым появлением $x$. 
\end{definition}

Пусть $B = \{b_1, \dots, b_{\beta}\}$ трёхрёберные блоки в порядке трансверсального обхода. $b_j = (e_{i_j}, e_{i_j + 1}, e_{i_j + 2})$. Обозначим $S_j = E(\gamma) \cap \{e_1, \dots, e_{i_j}\}$. 

\begin{conjecture}\label{conj:B-properties}
    $B$ имеет следующие свойства:
    \begin{itemize}
        \item $e_{i_j} \in E(\gamma)$.
        \item $e_{i_j + 1}, e_{i_j + 2} \not \in E(\gamma)$.
        \item Рёбра $e_{i_j + 1}, e_{i_j + 2}$ не являются петлями.
        \item Пусть ребро $e_{i_j + 2}$ соединяет вершины $x_j, y_j$. Существует полигональный путь $c_j \subset S_j$, такой что вершины $x_j, y_j$ являются его концами.
        \item Вершины $x_j, y_j$ являются вторыми появлениями для ребра $e_{i_j + 2}$.
    \end{itemize}    
\end{conjecture}


Докажем последний пункт. Раз путь $c_j$ проходит через вершины  $x_j, y_j$, число рёбер в $S_j$, инцидентных этим вершинам, ненулевое. Следовательно, эти вершины появляются до ребра $e_{i_j + 2}$.

Существование пути между вершинами $x, y$ в графе, индуцированном множеством рёбер $S_j$, будем обозначать $x \sim_j y$. Отношение $(\sim_j)$ является отношением эквивалентности, так как выражает связность в некотором графе. Для любых $j < k$ $S_j \subset S_k$. Следовательно, $\forall x, y \quad (x\sim_jy) \implies (x\sim_ky)$

\begin{lemma}\label{lemma:disjoint-c}
    Определённые выше пути $\{c_j\}_{j = 1}^\beta$ попарно не пересекаются по рёбрам.
\end{lemma}

\begin{proof}
    От противного. Пусть $e = (v_1, v_2)\in c_j \cap c_k$. Путь $c_j \subset c_k$, так как отношение $\sim_j$ слабее $\sim_k$. 
    Второе появление вершины $x_j$ является концевой вершиной в $c_k$, так как иначе ей было бы инцидентно два ребра из $S_k$. Это бы означало, что второму появлению $x_j$ инцидентно ребро из $E(\gamma)$, что не является таковым \ref{conj:B-properties}.
    Следовательно, второе появление $x_j$ лежит одновременно в блоках $b_j$ и $b_k$, которые не пересекаются. Противоречие.
\end{proof}

Обозначим через $m_i \in c_i$ ребро, которое появляется в трансверсальном пути первым. Все $c_i$ не пересекаются по рёбрам, а значит $m_i$ различны.

\begin{lemma}
    Обе вершины ребра $m_i$ являются первыми появлениями.
\end{lemma}

\begin{proof}
От противного. Пусть в ребре $m_i$ вершина $z$ появляется второй раз. Тогда $z \neq x_i, y_i$ (их вторые появления не входят в $c_i$), значит, существуют два ребра из $c_i$, инцидентные $z$: $ez_1$ (первое появление) и $ez_2 = m_i$ (второе). Однако $ez_1$ находится в пути раньше, чем $ez_2$. Противоречие.
\end{proof}

\begin{lemma}
    Если ребро $m_i$ не является первым в трансверсальном пути, оно следует после двухрёберного блока из $A$.
\end{lemma}

\begin{proof}
    От противного. Пусть ребро $m_i$ следует после блока $b_j \in B$. Тогда $y_j$ является вторым появлением в ребре $m_i$. Противоречие.
\end{proof}

\begin{corollary}\label{cor:a+1geqb}
    Для любого простого сборного графа без петель $\alpha + 1 \geq \beta$. 
\end{corollary}

\begin{proof}
    Мы инъективно сопоставили каждому блоку $b_i \in B$ ребро $m_i$, которое либо является первым в трансверсальном пути, либо следует после блока из $A$. Следовательно, $|B| = \beta \leq \alpha + 1$.
\end{proof}

\begin{theorem}\label{th:5k-3}
    Пусть $\Gamma$ простой сборный граф без петель, а $\gamma$ гамильтоново множество полигональных путей, построенное алгоритмом $\G$ для $\Gamma$. Тогда $|\Gamma| \geq 5 \cdot |\gamma| - 3$.
\end{theorem}

\begin{proof}
    Выразим все параметры через $\alpha$ и $\beta$

    \begin{itemize}
    \item Случай $e_{2n - 1} \not\in E(\gamma)$:
    \[
    \left\{
    \begin{aligned}
        & 2|\Gamma| - 1 = 2\alpha + 3\beta, \\
        & |E(\gamma)| = \alpha + \beta, \\
        & |\Gamma| = |\gamma| + |E(\gamma)| \\
        & \alpha + 1 \geq \beta 
    \end{aligned} 
    \right. 
    \implies{
        \left\{
        \begin{aligned}
            & 2|\Gamma| \geq 5\beta - 1, \\
            & 2\cdot|\gamma| = \beta + 1 \\
        \end{aligned} 
        \right.     
    }
    \implies |\Gamma| \geq 5|\gamma| - 3
    \]
    \item Случай $e_{2n - 1} \in E(\gamma)$:
    \[
    \left\{
    \begin{aligned}
        & 2|\Gamma| - 1 = 2\alpha + 3\beta + 1, \\
        & |E(\gamma)| = \alpha + \beta + 1, \\
        & |\Gamma| = |\gamma| + |E(\gamma)| \\
        & \alpha + 1 \geq \beta 
    \end{aligned} 
    \right.
    \implies{
        \left\{
        \begin{aligned}
            & 2|\Gamma| \geq 5\beta, \\
            & 2\cdot|\gamma| = \beta \\
        \end{aligned} 
        \right.
    } 
    \implies |\Gamma| \geq 5|\gamma|
    \]
\end{itemize}
\end{proof}

\begin{corollary}
    Пусть $\Gamma$ простой сборный граф без петель. Тогда $|\Gamma| \geq 5 \cdot \An(\Gamma) - 3$
\end{corollary}

Дальше мы увидим, что всякий граф $\hGamma$ с $|\hGamma| = 5 \cdot \An(\hGamma) - 3$ является $(1212)$-подстановкой какого-то другого графа (возможно, содержащего петли). Однако, в отличие от общего случая, алгоритм $\G$ плохо справляется с отделением таких графов от общей массы.

\begin{example}
    Пусть $\Gamma_1$ является сборным графом для слова $(1234125345)$. Рассмотрим следующий ряд:

    \[
    \Gamma_k = \underbrace{\Gamma_1 \circ \Gamma_1 \circ \dots \circ \Gamma_1}_{k - 1 \text{ раз}}\circ \Gamma_{(1212)} \quad  \forall k \geq 0. 
    \]

    где под $\Gamma_{(1212)}$ имеется в виду граф, соответствующий слову $(1212)$.

    \begin{itemize}
        \item В $\Gamma_k$ нет петель,
        \item $|\Gamma_k| = 5k - 3$,
        \item $\An(\Gamma_k) = 1$,
        \item $|\G(\Gamma_k)| = k$. 
    \end{itemize}
    
    Здесь $\G(\Gamma_k)$ обозначает гамильтоново множество полигональных путей, построенное алгоритмом $\G$.
\end{example}

\begin{proof}
    Из одного $\Gamma_1$ алгоритм $\G$ добавляет в $E(\gamma)$ следующие рёбра:

    \[
    \verb/( 1-2 3-4 1 2-5 3 4-5 )/
    \]

    Итого $|E(\gamma_{\G})| = 4 \cdot (k - 1) + 1, \quad |\gamma_{\G}| = |\Gamma_k| - |E(\gamma_{\G})| = k$. Оптимальный выбор рёбер выглядит так:

    \[
    \verb/(-1 2-3 4-1 2-5 3-4 5-)/
    \]

    Получаем $|E(\gamma_{opt})| = 5 \cdot (k - 1) + 1, \quad |\gamma_{opt}| = |\Gamma_k| - |E(\gamma_{opt})| = 1$.
    
    \todo{Сделать красивую картинку с работой двух алгоритмов}
\end{proof}


\begin{conjecture}\label{conj:G-properties}
    Пусть $\gamma$ гамильтоново множество полигональных путей, которое строит $\G$ для простого сборного графа без петель $\Gamma$, такое что $|\Gamma| = 5 \cdot |\gamma| - 3$. Тогда $\gamma$ имеет следующие свойства:

    \begin{enumerate}
        \item Последнее ребро трансверсального пути не содержится в $E(\gamma)$.
        \item $\alpha + 1 = \beta$.
        \item $2\cdot |\gamma| = \beta + 1$.
        \item После каждого двухрёберного блока $a \in A$ следует ребро $m_j$ для некоторого $j$.
    \end{enumerate}
\end{conjecture}

\begin{proof}
    Докажем каждый пункт по отдельности.
    \begin{enumerate}
        \item В противном случае $|\Gamma| \geq 5 \cdot |\gamma|$ \ref{th:5k-3}.
        \item Выводится из цепочки уравнений для случая $e_{2n - 1} \in E(\gamma)$ теоремы \ref{th:5k-3}.
        \item Следует из первого пункта и \ref{th:5k-3}.
        \item Каждое ребро $m_j$ либо является первым в трансверсальном пути, либо следует сразу после блока из $A$. Значит, как минимум после $|B| - 1 = \beta- 1$ блоков из $A$ следует ребро $m_j$. Однако $\beta - 1 = \alpha = |A|$. 
    \end{enumerate}
\end{proof}

Мы хорошо понимаем, что следует после блока из $A$. Следующим шагом будет сформировать схожее понимание для блоков из $B$. 
Блок из $A$ не может завершать трансверсальный путь (пункт 4 \ref{conj:G-properties}). 
Следовательно, последний блок $b_{\beta} \in B$ завершает трансверсальный путь. 
Вершину, что следует после блока $b_j \in B \quad \forall j\neq \beta$, будем обозначать $z_j$.

\begin{lemma}\label{lemma:z-properties}
    Для любой пары $j < \beta, k \leq \beta$ ни одно ребро из пути $c_k$ не инцидентно вершине $z_j$.
\end{lemma}

\begin{proof}
    От противного. Пусть существуют индексы $j,k$, такие что ребро $e \in c_k$ инцидентно $z_j$.
    Тогда $j < k$, так как не существует появления вершины из пути $c_k$, которое находится правее в трансверсальном пути, чем блок $b_k$. Рёбра $(y_j, z_j), e$ являются частью пути $c_k$. Следовательно, $c_j \subset c_k$. Это противоречит \ref{lemma:disjoint-c}.
\end{proof}

\begin{corollary}
    Пусть $\Gamma$ простой сборный граф без петель, такой что $|\Gamma| = 5\cdot |\G(\Gamma)| - 3$. Все пути $c_j$ имеют единичную длину.
\end{corollary}

\begin{proof}
    Рассмотрим произвольное ребро $e \in E(\gamma)$. Если оно является первым в трансверсальном пути, либо следует после блока из $A$, то оно является первым ребром пути $m_j \in c_j$. Иначе ребро следует после блока из $B$, а значит, не может быть частью какого-либо пути $c_j$ \ref{lemma:z-properties}. Следовательно, любой путь $c_j$ имеет единичную длину.
\end{proof}

\begin{lemma}
    Пусть $\Gamma$ будет простым сборным графом без петель, такой, что $|\Gamma| = 5\cdot |\G(\Gamma)| - 3$. Существует разбиение вершин $\{z_1,\ldots, z_{\beta - 1}\}$ на пары $(z_{p_k}, z_{q_k})$, такое что $z_{p_k} = z_{q_k} \quad \forall k$.
\end{lemma}

\begin{proof}
    Пусть рёбра $e_1, e_2$ инцидентны появлению $z_j$, отличному от того, что находится следом за $b_j$. Рёбра $e_1, e_2$ трансверсальны, $e_1$ находится раньше в трансверсальном пути, чем $e_2$.  Рассмотрим следующие случаи.

    \begin{itemize}
        \item $e_1, e_2 \in E(\gamma)$. Это нарушает полигональность путей из $\gamma$.
        \item $e_1, e_2 \not \in E(\gamma)$. Два трансверсальных ребра не находятся в $E(\gamma)$ тогда и только тогда, когда они находятся в одном и том же блоке из $B$. Значит, $z_j$ совпадает с вершиной $x_k$ для какого-то $b_k$. Это противоречит \ref{lemma:z-properties}, так как $x_k$ является концом пути положительной длины $c_k$.
        \item  $e_1 \not \in E(\gamma), e_2 \in E(\gamma)$. Так как $e_2$ отлично от всякого $m_k \in c_k$, оно должно следовать после какого-то блока $b_k \in B$. Однако это бы означало, что $z_j$ совпадает с $y_k$. Это противоречит \ref{lemma:z-properties}.
        \item  $e_1 \in E(\gamma), e_2 \not \in E(\gamma)$. Так как $e_1$ отличается от всякого $m_k \in c_k$, оно должно следовать после какого-то блока $b_k \in B$. Следовательно, вершина $z_j$ совпадает с вершиной $z_k$.
    \end{itemize}
\end{proof}

Мы вводим понятие \textit{паттерна} два-слова, которое упрощает его структуру до правильной скобочной последовательности. Дальше мы увидим, что этот объект сохраняет достаточно информации, чтобы с ним было удобно работать. 

\begin{definition}
    Слово $\mathcal{P}(w)$ в алфавите $\{1,2\}$, полученное из два-слова $w$ путём замены каждой буквы на номер её появления, будем называть \textit{паттерном} два-слова $w$.
\end{definition}

\begin{conjecture}
    Паттерн два-слова имеет следующие свойства:

    \begin{enumerate}
        \item $\mathcal{P}(w)$ является правильной скобочной последовательностью.
        \item Обозначим $\operatorname{Swap}: \{1, 2\}^* \xrightarrow{} \{1,2\}^*$ функцию, которая заменяет все буквы слова на противоположные. Тогда $\mathcal{P}(w^R) = \operatorname{Swap}(\mathcal{P}(w)^R)$. 
    \end{enumerate}
\end{conjecture}

\begin{proof} 
    Докажем каждый пункт по отдельности.
    \begin{enumerate}
        \item В два-слове одинаковое количество первых и вторых появлений. В префиксе два-слова количество первых появлений всегда не меньше количества вторых появлений.
        \item При чтении два-слова справа налево номер появления буквы меняется на противоположный.
    \end{enumerate}
\end{proof}

Мы убедились, что рёбра из множества $E(\gamma)$ больше связаны с блоками, предшествующими им, чем с теми, в которых они непосредственно находятся. Поэтому для описания сборных графов без петель с максимальным сборным числом будет удобно сдвинуть границы блоков на одно ребро вперёд.

\begin{theorem}\label{th:max-An-generic-struct}
    Пусть $\Gamma$ будет простым сборным графом без петель, таким, что $|\Gamma| = 5\cdot |\G(\Gamma)| - 3$. Для удобства обозначим $k = |\G(\Gamma)|$. Тогда выполняется следующий список свойств.

    \begin{enumerate}
        \item Слово $w$ является чередой из $2k - 1$ открывающих, $2k - 1$ закрывающих блоков и $k - 1$ главных букв, по два появления каждая.
        \item Открывающий блок состоит из двух первых появлений букв. Каждому открывающему блоку биективно соответствует закрывающий, который состоит из двух вторых появлений тех же букв.
        \item Слово $w$ начинается с открывающего блока и заканчивается закрывающим.
        \item После каждого закрывающего блока, кроме последнего, следует появление одной из главных вершин.
    \end{enumerate}
\end{theorem}

\begin{proof}
    Заметим, что $\alpha + 1 = \beta = 2k - 1$. Открывающие блоки соответствуют рёбрам $\{m_j\}$, а закрывающие --- крайним рёбрам блоков из $B$. Главные буквы соответствуют вершинам $z_i$.

    \begin{enumerate}
        \item Так как $\alpha + 1 = \beta$, слово $w$ действительно начинается с ребра $m_j$ для какого-то $j$.
        \item Все $x_j, y_j$ являются концами однорёберного пути $c_j$. Следовательно, $x_j, y_j$ действительно являются вторыми появлениями вершин из $m_j$.
        \item Вершины $z_j$ следуют после блока $b_j$ при $j \neq \beta$. Подслово $\{z_j\}$ действительно является два-словом размера $k - 1$.
    \end{enumerate}
\end{proof}

Так как транспонирование два-слова $w$ не меняет его сборное число, логично будет изучить поведение жадного алгоритма в работе со словом $w^R$. Действительно, если сборное число достигает максимального значения $\An(w) = \frac{1}{5}(|w| + 3)$, то $|\G(w)| = |\G(w^R)| = \An(w)$.

\begin{definition}
    Пусть $\Gamma$ --- простой сборный граф. \textit{Симметричным жадным алгоритмом}  $\G_s$ будем называть итеративный алгоритм, который строит гамильтоновы множества полигональных путей $\G(\Gamma), \G(\Gamma^R)$ и выводит то, которое имеет наименьший размер.
\end{definition}


\begin{corollary}
    Если $|\Gamma_w| = 5\cdot |\G_s(\Gamma_w)| - 3$, то
    \begin{enumerate}
        \item $\pattern(w)$ соответствует регулярному выражению \verb/^11(11|221|222)/$^*$\verb/22$/
        \item $\pattern(w)$ соответствует регулярному выражению \verb/^11(22|111|211)/$^*$\verb/22$/
    \end{enumerate}
\end{corollary}

\begin{proof}
    Открывающий блок вносит вклад \verb/11/ в паттерн, а закрывающий блок с его главной вершиной \verb/221/ или \verb/222/. Сборное число не меняется при транспонировании, так что $\pattern(w^R)$ тоже соответствует регулярному выражению \verb/^11(11|221|222)/$^*$\verb/22$/. Транспонируя обе части выражения, получаем, что $\pattern(w)$ соответствует регулярному выражению \verb/^11(22|111|211)/$^*$\verb/22$/.
\end{proof}

\begin{conjecture}
    Пересечение языков, соответствующих выражениям \verb/^(11|221|222)/$^*$\verb/$/ и \verb/^(22|111|211)/$^*$\verb/$/, порождается регулярным выражением \verb/^(22?11|111111|222222|2222211111)/$^*$\verb/$/
\end{conjecture}

\begin{proof}
    Доказывается несложным перебором.
\end{proof}

\begin{figure}[htbp]
    \centering    \includegraphics[width=0.56\textwidth]{img/ff.pdf}
    \caption{Префиксное дерево слов 111111, 22111, 22211, 2222211111, 222222}
    \label{fig}
\end{figure}

Не все блоки из полученного списка хорошо сочетаются с остальными свойствами конструкции из теоремы \ref{th:max-An-generic-struct}. Рассмотрим количество открывающих и закрывающих блоков каждого из них.

\begin{itemize}
    \item \verb/22111/ и \verb/22211/ состоят из одного открывающего и одного закрывающего блока, как при чтении слева направо, так и при чтении справа налево.
    \item \verb/2222211111/ состоит из двух открывающих и двух закрывающих блоков, как при чтении слева направо, так и при чтении справа налево.
    \item \verb/222222/ состоит из двух закрывающих блоков, а его реверс \verb/111111/ состоит из трёх открывающих блоков.
\end{itemize}

Итого, у блоков \verb/22111/, \verb/22211/ и \verb/2222211111/ суммарно при чтении слева направо и справа налево одинаковое количество открывающих и закрывающих блоков, в то время как блоки \verb/222222/ и \verb/111111/ имеют на один открывающий блок больше.

\begin{conjecture}
    Если $|\Gamma_w| = 5\cdot |\G_s(\Gamma_w)| - 3$, то в $\pattern(w)$ не встречаются блоки \verb/222222/ и \verb/111111/. Другими словами, $\pattern(w)$ соответствует регулярному выражению \verb/^11(22?11|2222211111)/$^*$\verb/22$/
\end{conjecture}

\begin{proof}
    От противного. Если блоки \verb/222222/ и \verb/111111/ содержатся в $\pattern(w)$, то либо при чтении слева направо, либо при чтении справа налево у слова $w$ будет больше открывающих блоков, чем закрывающих. Однако в описанной конструкции открывающих и закрывающих блоков одинаковое количество. Противоречие.
\end{proof}

У каждого из трёх оставшихся блоков одинаковое количество открывающих и закрывающих блоков. 

\begin{conjecture}
    Если $|\Gamma_w| = 5\cdot |\G_s(\Gamma_w)| - 3$, то в $\pattern(w)$ не встречаются блоки \verb/2222211111/. Другими словами, $\pattern(w)$ соответствует регулярному выражению \verb/^11(22?11)/$^*$\verb/22$/
\end{conjecture}

\begin{proof}
    От противного. Пусть $b$ блок вида \verb/2222211111/. В префиксе $\pattern(w)$ перед $b$ на один больше открывающих блоков, чем закрывающих. Так как $b$ начинается с двух закрывающих блоков, существует префикс слова $w$, содержащий больше закрывающих блоков, чем открывающих. Это противоречит пункту 2 теоремы \ref{th:max-An-generic-struct}.
\end{proof}

Регулярное выражение \verb/^11(22?11)/$^*$\verb/22$/ можно записать как \verb/^(1122?)/$^*$\verb/1122$/. Так как в слове вида \verb/^(1122?)/$^*$\verb/1122$/ открывающие и закрывающие блоки чередуются, между парой соседних главных вершин находится копия слова $1212$. 
Это в точности совпадает с определением $(1212)$-подстановки $\mathcal{S}_{(1212)}$. 
С другой стороны, мы знаем, что все $1212$-подстановки имеют максимальное возможное сборное число \ref{cor:pretzelsat5k-3}.

\begin{corollary}
    Простой сборный граф без петель $\Gamma$ имеет максимальное сборное число $5\cdot \An(\Gamma) - 3 = |\Gamma|$ тогда и только тогда, когда существует простой сборный граф $\bar{\Gamma}$ (возможно, содержащий петли), такой, что $\Gamma = \mathcal{S}_{(1212)}(\bar{\Gamma})$
\end{corollary}

Теперь перейдём к изучению случая графов без петель с максимальным возможным сборным числом и $|\Gamma| \equiv \D - 3 \mod{5}$, при $0 < \D < 5$.

\begin{definition}
    Для каждого целого числа $\D, 0 \leq \D < 5$ обозначим через $\OPT_\D$ семейство простых сборных графов без петель $\Gamma$, таких, что $\Gamma$ и $\Gamma^R$ имеют следующие свойства:

    \begin{enumerate}
        \item Сборное число принимает максимальное значение: $|\Gamma| = \D + 5\cdot \An(\Gamma) - 3$.
        \item Последние три ребра трансверсального пути формируют блок из семейства трёхрёберных блоков, заданного алгоритмом $\G$.
        \item Первое ребро в трансверсальном пути является $m_i$ для некоторого пути $c_i$, где множества $\{m\}_i$ и $\{c\}_i$ заданы алгоритмом $\G$.  
    \end{enumerate}
\end{definition}

При $\D > 0$ уже не все простые сборные графы без петель с максимальным сборным числом находятся в $\OPT_\D$. Чтобы упростить дальнейшее рассуждение, мы приведём их к нужному виду с помощью композиции с аддитивными словами.

\begin{conjecture}
    Пусть сборный граф без петель $\Gamma$ имеет максимальное сборное число $|\Gamma| = \D +5 \cdot \An(\Gamma) - 3$ для некоторого $0 < \Delta < 5$.
    Пусть $\hGamma = \Gamma_{1212} \circ \mathcal{S}_{1212}^\circ(1212) \circ\Gamma \circ \mathcal{S}_{1212}^\circ(1212) \circ \Gamma_{1212}$. Тогда $\hGamma \in \OPT_\D$.
\end{conjecture}

\begin{proof}
    Слово $w = \Gamma_{1212} \circ \mathcal{S}_{1212}^\circ(1212)$ является аддитивным слева, а слово $w^R = \mathcal{S}_{1212}^\circ(1212) \circ \Gamma_{1212}$ аддитивным справа. Следовательно, $\An(\hGamma) = \An(\Gamma_w \circ \Gamma \circ \Gamma_{w^R}) = \An(\Gamma_w) + \An(\Gamma) + \An(\Gamma_{w^R}) = 4 + \An(\Gamma)$.
    
    \begin{enumerate}
        \item  Проверяем равенство: $|\hGamma| = 20 + |\Gamma| = \D+ 5 \cdot \An(\hGamma) - 3$
        \item  То, какие рёбра подграфа $w^R = \mathcal{S}_{1212}^\circ(1212) \circ \Gamma_{1212}$ будут добавлены алгоритмом $\G$ в $E(\gamma)$, зависит только от того, находится ли ребро, соединяющее его с остальными вершинами, в $E(\gamma)$ или нет.

        \[
            \verb| 1-2 3-2 3-4 5-6 5 6-1 7-8 7 8-4 9-10 9 10|
        \]
        \[
            \verb|-1 2-3 2 3-4 5-6 5 6-1 7-8 7 8-4 9-10 9 10|
        \]

        В обоих случаях трансверсальный путь замыкает трёхрёберный блок из $B$.
        
        \item  Граф $\hGamma$ начинается с $\Gamma_{1212}$. Его первые три ребра формируют блок из $B$. Следовательно, первое ребро трансверсального пути является началом пути $c_1$.
    \end{enumerate}

    Утверждения для $\hGamma^R$ следуют из того, что $\hGamma^R = \Gamma_w \circ \Gamma^R \circ \Gamma_{w^R}$.
\end{proof}

 Любой простой сборный граф без петель $\Gamma$, для которого существует $\hGamma \in \OPT_\D$ такой, что $\hGamma = \Gamma_{1212} \circ \mathcal{S}_{1212}^\circ(1212) \circ\Gamma \circ \mathcal{S}_{1212}^\circ(1212) \circ \Gamma_{1212}$, по аддитивности имеет максимальное сборное число. С другой стороны, для всех $\Gamma$ с максимальным сборным числом $\hGamma \in \OPT_\D$. Поэтому, чтобы изучить семейство графов с максимальным сборным числом, достаточно изучить $\OPT_\D$.

 \begin{conjecture}
     Пусть простой сборный граф без петель $\Gamma \in \OPT_\D$ для некоторого $0 \leq \D \leq 4$ и $\alpha, \beta$ --- мощности семейств двухрёберных и трёхрёберных блоков, задаваемых алгоритмом $\G$,  соответственно. Тогда $\alpha + 1 = \beta + \D$.
 \end{conjecture}

\begin{proof}
    Для удобства обозначим $k = \An(\Gamma)$. Воспользуемся утверждениями \ref{conj:no-loops-G-properties}, \ref{cor:a+1geqb}:

    \[
        \left\{
        \begin{aligned}
            & 2\cdot|\Gamma| - 1 = 2\alpha + 3\beta, \\
            & |\Gamma| - |\gamma| = \alpha + \beta, \\
            & |\Gamma| = 5\cdot k - 3 + \D, \\
            & \alpha + 1 \geq \beta 
        \end{aligned} 
        \right.
        \implies{} 
        \left\{
        \begin{aligned}
            & 10 \cdot k - 7 + 2\D = 2\alpha + 3\beta, \\
            & 2|\gamma| - 1 = \beta, \\
            & \alpha + 1 \geq \beta 
        \end{aligned} 
        \right.
    \]

    Рассмотрим первое равенство по модулю 5.
    \[
        \left\{
        \begin{aligned}
            & 2\beta + 2\D \equiv 2\alpha  + 2 \mod{5}, \\
            & \alpha + 1 \geq \beta 
        \end{aligned} 
        \right.
        \implies{} 
        \left\{
        \begin{aligned}
            & \alpha  + 1\equiv \beta + \D  \mod{5}, \\
            & \alpha + 1 \geq \beta 
        \end{aligned} 
        \right.
        \implies{\alpha + 1 \geq \beta + \D}
    \]

    Применим усиленное неравенство.

    \[
        \left\{
        \begin{aligned}
            & 5 \cdot (2k - 1) + 2\D \geq 5\beta + 2\D,\\
            & 2\cdot|\gamma| - 1 = \beta, \\
            & \alpha + 1 \geq \beta + \D 
        \end{aligned} 
        \right.
        \implies{k \geq |\gamma|}
    \]

    Сборное число $k = \An(\Gamma)$ по определению не больше размера любого гамильтонова множества полигональных путей $\gamma$, а значит $k = |\gamma|$. Следовательно, неравенство превращается в равенство: $\alpha + 1 = \beta + \D$
    
\end{proof}




\bibliographystyle{unsrt}
\bibliography{refs}
\end{document}